% Part: first-order-logic
% Chapter: completeness
% Section: completeness-theorem

\documentclass[../../../include/open-logic-section]{subfiles}

\begin{document}

\iftag{FOL}
      {\olfileid{fol}{com}{cth}}
      {\olfileid{pl}{com}{cth}}

\olsection{संपूर्णता प्रमेय}

\begin{explain}
आपल्या निष्पत्ती एकत्र करू: आपल्याला संपूर्णता प्रमेय मिळते.
\end{explain}

\begin{thm}[संपूर्णता प्रमेय]
\ollabel{thm:completeness}
$\Gamma$ हा !!{sentence}s चा संच असू द्या. $\Gamma$ सुसंगत असेल, तर
तो पूर्ततायोग्य आहे.
\end{thm}

\begin{proof}
$\Gamma$ सुसंगत आहे असे समजा.\iftag{FOL}{
  \olref[hen]{lem:henkin} नुसार असा संतृप्त सुसंगत संच
  $\Gamma' \supseteq \Gamma$ असतो.}{} \olref[lin]{lem:lindenbaum} नुसार
असा $\Gamma^* \supseteq \iftag{FOL}{\Gamma'}{\Gamma}$ असतो की तो
सुसंगत आणि !!{complete} असतो.
\iftag{FOL}{
  $\Gamma' \subseteq \Gamma^*$ असल्यामुळे प्रत्येक
  !!{formula}~$!A(x)$ साठी~$\Gamma^*$ मध्ये पुढील रूपाचे !!a{sentence}
  असते:
  \iftag{prvEx}
        {$\lexists[x][!A(x)] \lif !A(c)$}
        {$\lnot\lforall[x][!A(x)] \lif \lnot !A(c)$}.
  म्हणून~$\Gamma^*$ संतृप्त आहे. $\Gamma$ मध्ये~$\eq$ येत नसेल, तर}{या}
\olref[mod]{lem:truth} नुसार
$\iftag{FOL}{\Sat{M(\Gamma^*)}}{\pSat{v(\Gamma^*)}}{!A}$ तेव्हा आणि केवळ तेव्हाच, जेव्हा $!A \in
\Gamma^*$. यावरून विशेषतः प्रत्येक $!A \in \Gamma$ साठी
$\iftag{FOL}{\Sat{M(\Gamma^*)}}{\pSat{v(\Gamma^*)}}{!A}$; म्हणून
$\Gamma$ पूर्ततायोग्य आहे.\iftag{FOL}{
  $\Gamma$ मध्ये~$\eq$ येत असेल, तर \olref[ide]{lem:truth} नुसार
  सर्व !!{sentence}s~$!A$ यांच्यासाठी
  $\Sat{\equivclass{M}{\approx}}{!A}$ तेव्हा आणि केवळ तेव्हाच, जेव्हा
  $!A \in \Gamma^*$. विशेषतः, प्रत्येक $!A \in \Gamma$ साठी
  $\Sat{\equivclass{M}{\approx}}{!A}$; म्हणून $\Gamma$ पूर्ततायोग्य आहे.}{}
\end{proof}

\begin{cor}[संपूर्णता प्रमेय, दुसरी आवृत्ती]
\ollabel{cor:completeness}
प्रत्येक $\Gamma$ आणि प्रत्येक !!{sentence}s~$!A$ साठी: जर
$\Gamma \Entails !A$, तर $\Gamma \Proves !A$.
\end{cor}

\begin{proof}
लक्षात घ्या की \olref{cor:completeness} आणि
\olref{thm:completeness} यांमधील $\Gamma$ वर सर्ववाची संख्यकीकरण केलेले
आहे. आपला गोंधळ होऊ नये म्हणून \olref{thm:completeness} हे वेगळा चर
वापरून पुन्हा मांडू: !!{sentence}s चा कोणताही संच~$\Delta$ दिला असता,
जर $\Delta$ सुसंगत असेल, तर तो पूर्ततायोग्य आहे. प्रतिपरिवर्तनाने,
$\Delta$ पूर्ततायोग्य नसेल, तर $\Delta$ विसंगत आहे. अनुप्रमेय सिद्ध
करण्यासाठी याचा वापर करू.

$\Gamma \Entails !A$ असे समजा. मग \olref[syn][sem]{prop:entails-unsat}
नुसार $\Gamma \cup \{\lnot !A\}$ अपूर्ततायोग्य आहे. $\Gamma \cup
\{\lnot !A\}$ हाच आपला $\Delta$ घेतल्यास,
\olref{thm:completeness} च्या मागील आवृत्तीनुसार $\Gamma \cup
\{\lnot !A\}$ विसंगत आहे. पुढील संदर्भांनुसार,
\iftag{FOL}{%
  \tagrefs{prfAX/{fol:axd:prv:prop:prov-incons},
    prfSC/{fol:seq:prv:prop:prov-incons},
    prfND/{fol:ntd:prv:prop:prov-incons},
    prfTab/{fol:tab:prv:prop:prov-incons}}}{%
  \tagrefs{prfAX/{pl:axd:prv:prop:prov-incons},
    prfSC/{pl:seq:prv:prop:prov-incons},
    prfND/{pl:ntd:prv:prop:prov-incons},
    prfTab/{pl:tab:prv:prop:prov-incons}}},
$\Gamma \Proves !A$.
\end{proof}

\tagprob{FOL}
\begin{prob}
\olref[fol][com][cth]{cor:completeness} वापरून
\olref[fol][com][cth]{thm:completeness} सिद्ध करा आणि अशा रीतीने
संपूर्णता प्रमेयाच्या दोन्ही मांडण्या समतुल्य आहेत हे दाखवा.
\end{prob}
\tagendprob

\tagprob{notFOL}
\begin{prob}
\olref[pl][com][cth]{cor:completeness} वापरून
\olref[pl][com][cth]{thm:completeness} सिद्ध करा आणि अशा रीतीने
संपूर्णता प्रमेयाच्या दोन्ही मांडण्या समतुल्य आहेत हे दाखवा.
\end{prob}
\tagendprob

\tagprob{FOL}
\begin{prob}
एखादी !!a{derivation} प्रणाली संपूर्ण असण्यासाठी तिचे नियम इतके समर्थ
असले पाहिजेत की प्रत्येक अपूर्ततायोग्य संच विसंगत असल्याचे ते सिद्ध करू
शकतील. संपूर्णता सिद्ध करण्यासाठी !!{derivation} चे कोणते नियम आवश्यक
होते? यांपैकी कोणताही नियम सिद्धतेत कोठेच वापरलेला नाही का? या प्रश्नांची
उत्तरे देण्यासाठी, \olref[fol][com][cth]{thm:completeness} च्या सिद्धतेपर्यंत
नेणाऱ्या कोणत्या निष्पत्तींमध्ये !!{derivation} चे कोणते नियम वापरले, हे
दाखवणारी यादी किंवा आकृती तयार करा. या सिद्धतांमधील नियमांच्या कोणत्याही
अध्याहृत वापरांची नोंद अवश्य करा.
\end{prob}
\tagendprob

\tagprob{notFOL}
\begin{prob}
एखादी !!a{derivation} प्रणाली संपूर्ण असण्यासाठी तिचे नियम इतके समर्थ
असले पाहिजेत की प्रत्येक अपूर्ततायोग्य संच विसंगत असल्याचे ते सिद्ध करू
शकतील. संपूर्णता सिद्ध करण्यासाठी !!{derivation} चे कोणते नियम आवश्यक
होते? यांपैकी कोणताही नियम सिद्धतेत कोठेच वापरलेला नाही का? या प्रश्नांची
उत्तरे देण्यासाठी, \olref[pl][com][cth]{thm:completeness} च्या सिद्धतेपर्यंत
नेणाऱ्या कोणत्या निष्पत्तींमध्ये !!{derivation} चे कोणते नियम वापरले, हे
दाखवणारी यादी किंवा आकृती तयार करा. या सिद्धतांमधील नियमांच्या कोणत्याही
अध्याहृत वापरांची नोंद अवश्य करा.
\end{prob}
\tagendprob

\end{document}
